\documentclass{article}
\usepackage{/home/user1/code/tex/sty}
\usepackage{amsfonts}
\usepackage{amsmath}
\usepackage{mathtools}
\usepackage{geometry}
\usepackage{graphicx}
\usepackage{hyperref}
\usepackage{floatrow}
\newcommand{\fig}[3]{
\begin{figure}[H]
\begin{center}
\includegraphics[height=#1]{#2}
\caption{#3}
\label{#2}
\end{center}
\end{figure}}
\setlength{\parindent}{0pt}
\geometry{top=1in,bottom=1in,left=1in,right=1in}
\begin{document}
\begin{center}
\textbf{\Large{Problem Set 9}}\\~\\
\begin{tabular}{l l}
Student&Agustin Romero\\
Email&ar1@stanford.edu\\
Instructor&Lernik Asserain\\
Course&MATH 131P\\
Institution&Stanford University\\
Date&\today\\
\end{tabular}
\end{center}
\section*{Preface}
\e{\lab{aosudn}\int x^{n+1}J_n(x)dx = x^{n+1}J_{n+1}(x)+C}
\e{\lab{ansd8n}\int_0^a J_n(\fr{\alpha_j}{a}x)^2 x dx = \fr{a^2}{2}J_{n+1}(\alpha_j)^2}
\e{\lab{mfu1z}\int_0^a J_n(\fr{\alpha_j}{a}x)J_n(\fr{\alpha_k}{a}x)xdx=0 \quad j\neq k}
\section*{Problem 1.a.}
\e{x=r\cos(\theta)\quad y=r\sin(\theta)}
\e{u(x,y)&=\fr{x}{x^2+y^2}=\fr{r\cos(\theta)}{(r\cos(\theta))^2+(r\sin(\theta))^2}=\fr{r\cos(\theta)}{r^2(\cos(\theta)^2+\sin(\theta)^2)}\\
&=\boxed{\fr{\cos(\theta)}{r}}}
The Laplace equation is
\n{laplace}{(\fr{\p^2}{\p r^2} + \fr{1}{r}\fr{\p}{\p r} + \fr{1}{r^2}\fr{\p^2}{\p \theta^2})u}
\e{\laplace &= 0\\
2r^{-3}\cos(\theta)-r^{-3}\cos(\theta)-r^{-3}\cos(\theta)&= \boxed{0}}
\fbox{$u$ satisfies the Laplace equation.}
\section*{Problem 1.b.}
\e{\ln(x^2+y^2)=\ln(r^2)}
\e{\fr{\p}{\p r} \ln(r^2) = \fr{1}{r^2} 2r = 2r^{-1}}
\e{\fr{\p^2}{\p r^2} = -2r^{-2}}
The Laplace equation is
\e{\laplace &= 0\\
-2r^{-2}+2r^{-2}+0 = \boxed{0}}
\fbox{$u$ satisfies the Laplace equation.}
\section*{Problem 2.a.}
This is a two dimensional heat equation in polar coordinates, independent of $\theta$. Find $A_j$, then $u(t,r)$.
\n{aj}{\alpha_j}
\e{A_j &= \fr{2}{a^2 J_1(\aj)^2} \int_0^a f(r) J_0(\aj r)r dr\\
&= \fr{200}{J_1(\aj)^2} \int_0^1 J_0(\aj r) r dr
\intertext{Use $u=\aj r$, $du=\aj dr$ and then \eref{aosudn}.}
&= \fr{200}{\aj^2 J_1(\aj)^2} (u J_1(u))_0^{\aj}\\
&= \boxed{\fr{200}{\aj J_1(\aj)}}}
Find $u(t,r)$.
\e{u(t,r)&= \sum_{j=1}^\infty A_j e^{-(\aj / a)^2 t} J_0(\aj a^{-1} r)\\
&= \sum_{j=1}^\infty A_j e^{- \aj^2 t} J_0(\aj r)\\
&= \boxed{\sum_{j=1}^\infty \fr{200}{\aj J_1(\aj)} e^{-\aj^2 t} J_0(\aj r)}}
\section*{Problem 2.b.}
\e{J_0(r)=\fr{1}{2\pi} \int_{-\pi}^\pi e^{2x \sin(t)} dt}
\e{|J_0(r)|=1=J_0(0)}
\fbox{$J_0(r)$ is maximized at $r=0$, so the hottest temperature at any time is at $r=0$.} Intuitively, the temperature is highest at lower $t$, and farther away from the boundary, so \fbox{lower $r$}.
\section*{Problem 3.a.}
This is a two dimensional wave equation in polar coordinates, independent of $\theta$, with vanishing Dirichlet boundary conditions. Find $A_j$, $B_j$, $u(t,r)$.
\n{eindimensionaj}{\fr{2}{a^2 J_1(\aj)^2}\int_0^a f(r) J_0(\fr{\aj}{a}r)r dr}
\n{eindimensionbj}{\fr{2}{a \aj J_1(\aj)^2}\int_0^a g(r) J_0(\fr{\aj}{a}r)r dr}
\n{sjeu}{\sum_{j=1}^\infty}
\n{eindimensionu}{\sjeu (A_j \cos(\fr{\aj}{a} t)+B_j \sin(\fr{\aj}{a} t))J_0(\fr{\aj}{a} r)}
\e{A_j&= \eindimensionaj \\
&= \fr{2}{4 J_1(\aj)^2} \int_0^a 0 dr\\
&= \boxed{0}}
Find $B_j$.
\e{B_j&= \eindimensionbj \\
&= \fr{2}{2 J_1(\aj)^2} \int_0^2 J_0(\fr{\aj}{2} r)r dr
\intertext{Use u-substitution, $u=\fr{\aj}{2} r$, $du=\fr{\aj}{2} dr$.}
&= \fr{1}{\aj J_1(\aj)^2} (\fr{2}{\aj})^2 \int_0^{\aj} J_0(u) u du
\intertext{Use \eref{aosudn}.}
&= \boxed{\fr{4}{\aj^2 J_1(\aj)}}}
Find $u(t,r)$.
\e{u(t,r)&=\eindimensionu \\
&=\boxed{\sum_{j=1}^\infty \fr{4}{\aj^2 J_1(\aj)} \sin(\fr{aj}{2}t)J_0(\fr{\aj}{2}r)}}
\section*{Problem 3.b.}
This is a two dimensional wave equation in polar coordinates, independent of $\theta$, with vanishing Dirichlet boundary conditions. Find $A_j$, $B_j$, $u(t,r)$.
\e{A_j=\eindimensionaj=\boxed{0}}
\e{B_j&=\eindimensionbj\\
&=\fr{2}{\aj J_1(\aj)^2}\int_0^a J_0(\fr{\alpha_3}{a} r)J_0(\fr{\aj}{a} r)rdr
\intertext{Use \eref{mfu1z}.}
&=\falls{0&\quad j\neq 3\\
\fr{2}{\aj J_1(\aj)^2}\int_0^a J_0(\fr{\alpha_3}{a}r)^2 r dr &\quad j=3}
\intertext{Use \eref{ansd8n}.}
&= \falls{0 &\quad j\neq 3\\
\fr{2}{\aj J_1(\aj)^2} \fr{a^2}{2} J_1(\aj)^2 &\quad j=3}\\
&=\boxed{\falls{0 &\quad j\neq 3\\
\fr{1}{\aj} &\quad j=3}}}
\e{u(t,r)&=\eindimensionu\\
&=\boxed{\fr{1}{\alpha_3} \sin(\alpha_3t )J_0(\alpha_3 r)}}
\section*{Problem 4}
This is a two dimensional wave equation in polar coordinates, independent of $\theta$, with vanishing Dirichlet boundary conditions. Find $A_j$, $B_j$, and $u(t,r)$.
\e{\fr{\p}{\p \theta} u=0 \quad c=10 \quad a=1 \quad f=1-r^2 \quad g=1}
\e{A_j&=\eindimensionaj\\
&=\fr{2}{J_1(\aj)^2}\int_0^1 (1-r^2)J_0(\aj r)r dr
\intertext{Expand $1-r^2$ as suggested, making sure to use a different index $k$.}
&= \fr{2}{J_1(\aj)^2} \sum_{k=1}^\infty (\fr{8}{\alpha_k^3 J_1(\alpha_j)} \int_0^a J_0(\fr{\alpha_k}{a}r)J_0(\fr{\alpha_j}{a}r)r dr)
\intertext{Use \eref{mfu1z} to find $j=k$, and then integrate $J^2$.}
&= \fr{16}{\aj^3 J_1(\aj)^3} \fr{a^2}{2}J_1(\aj)^2\\
&= \boxed{\fr{8}{\aj^3 J_1(\aj)}}}
Find $B_j$.
\e{B_j&=\eindimensionbj\\
&= \fr{2}{c\aj J_1(\aj)^2} \int_0^1 J_0(\aj r)r dr\\
&= \fr{1}{5 \aj^3 J_1(\aj)^2} \int_0^{\aj} J_0(u)u du\\
&= \fr{1}{5 \aj^3 J_1(\aj)^2} (J_1(u)u)_0^{\aj}\\
&= \boxed{\fr{1}{5 \aj^2 J_1(\aj)}}}
Find $u(t,r)$.
\e{u(t,r)&=\eindimensionu\\
&=\boxed{\sum_{j=1}^\infty (\fr{8}{\aj^3 J_1(\aj)} \cos(10\aj t)+\fr{1}{5\aj^2 J_1(\aj)} \sin(10\aj t)) J_0(\aj r)}}
\section*{Problem 5}
This is a two dimensional Laplace equation in polar coordinates with Dirichlet boundaries. Find $A_0$, $A_n$, $B_n$, and $u(r,\theta)$.
\n{lapanull}{\fr{1}{2\pi}\int_0^{2\pi} f(\theta)d\theta}
\n{lapan}{\fr{1}{\pi}\int_0^{2\pi}f(\theta)\cos(n\theta)d\theta}
\n{lapbn}{\fr{1}{\pi}\int_0^{2\pi}f(\theta)\sin(n\theta)d\theta}
\n{lapu}{A_0+\sum_{n=1}^\infty (\fr{r}{a})^2(A_n \cos(n\theta) +B_n\sin(n\theta))}
\e{A_0&= \lapanull\\
&= \fr{1}{2\pi}(\pi \int_0^\pi d\theta - \int_0^\pi \theta d\theta)\\
&= \fr{1}{2\pi}(\pi^2 - \fr{1}{2}\pi^2)\\
&= \boxed{\fr{\pi}{4}}}
Find $A_n$.
\e{A_n&=\lapan\\
&= \fr{1}{\pi}\int_0^\pi (\pi-\theta)\cos(n\theta)d\theta\\
&= \fr{1}{\pi}((\pi-\theta)\sin(n\theta)-\fr{1}{n^2}\cos(n\theta))_0^\pi\\
&= \fr{1}{\pi}(-\fr{1}{n^2}\cos(n\theta)+\fr{1}{n^2})\\
&= \fr{1}{\pi n^2}(1-(-1)^n)\\
&= \boxed{\falls{0&\quad \textrm{n odd}\\
\fr{2}{\pi n^2} &\quad \textrm{n even}}}}
Find $B_n$.
\e{B_n&=\lapbn\\
&= \fr{1}{\pi}\int_0^\pi (\pi-\theta)\sin(n\theta)d\theta\\
&= \fr{1}{\pi}(-\fr{1}{n}(\pi - \theta)\cos(n\theta)-\fr{1}{n^2}\sin(n\theta))_0^\pi\\
&= \fr{1}{\pi}(\fr{1}{n}\pi)\\
&= \boxed{\fr{1}{n}}}
Find $u(r,\theta)$.
\e{u(r,\theta)&=\lapu\\
&=\boxed{\fr{\pi}{4}+\sum_{k=0}^\infty \fr{2r^{2k+1}}{\pi(2k+1)^2}\cos((2k+1)\theta)+\sum_{n=1}^\infty \fr{r^n}{n}\sin(n\theta)}}
\section*{Problem 6.a.}
Write the two dimensional Laplace equation in polar coordinates, simplify.
\e{\fr{\p}{\p r}u(r,\theta)=0}
\e{\nabla^2 u &= 0\\
\laplace &= 0\\
\fr{1}{r^2}\fr{\p^2}{\p \theta^2} u &= 0\\
\fr{\p^2}{\p \theta^2} u &= 0}
This is a second order ODE. Solve using the characteristic polynomial.
\e{\rho^2+c_1\rho+c_2&=0\\
\rho^2&= 0\\
\rho&=0}
Repeated roots. The solution is
\e{u(r,\theta)&=u(\theta)\\
&= a \theta e^{\rho \theta} + b e^{\rho \theta}\\
&= \boxed{a \theta + b}}
where $a$,$b\in \mathbb{R}$.
\section*{Problem 6.b.}
Apply boundary conditions.
\e{u(r,0)=T_1 \Rightarrow \boxed{b = T_1}}
\e{u(r,\theta_0)=T_1+a\theta_0 = T_2 \Rightarrow \boxed{a = \fr{1}{\theta_0}(T_2-T_1)}}
The solution for the Laplace equation in the wedge is
\e{\boxed{u(\theta)=T_1+\fr{1}{\theta_0}(T_2-T_1)}}
\section*{Problem 7.a.}
Write the two dimensional Laplace equation in polar coordinates, simplify.
\e{\fr{\p}{\p \theta}u(r,\theta)=0}
\e{\nabla^2 u &= 0\\
\laplace &= 0\\
(\fr{\p^2}{\p r^2}+\fr{1}{r}\fr{\p}{\p r})u&=0\\
r^2 u'' + ru' &= 0}
This is a Cauchy-Euler equation. Solve the Indicial equation and find $u(r)$.
\e{r^2 u'' + au' + b&=0 \Rightarrow a=1 \quad b=0}
\e{\textrm{Indicial}\Rightarrow \rho^2 + (a-1)\rho +b &=0\\
\rho^2 &= 0\\
\rho &= 0}
The solution is
\e{u(r)&=a r^\rho \ln(r) + b r^\rho\\
&= \boxed{a \ln (r) + b}}
\section*{Problem 7.b.}
Apply boundary conditions.
\e{u(1)=a\ln(1)+b=\boxed{b=2}}
\e{u(\fr{1}{2})&=2+a\ln(\fr{1}{2})=1\\
a\ln(\fr{1}{2})&=-1\\
a&=-\fr{1}{\ln(\fr{1}{2})}=\boxed{\fr{1}{\ln(2)}}}
\section*{Problem 8}
This is a two dimensional Laplace equation in polar coordinates with mixed vanishing Dirichlet and Neumann boundary conditions. Find $B_n$, and the time-independent solution $u(r,\theta)$.
\e{u(r,\theta)&=\sum_{n=1}^\infty B_n r^{n \pi \theta_0^{-1}} \sin(n \pi \theta_0^{-1} \theta)\\
&= \sum_{n=1}^\infty B_n r^{4n} \sin(4n\theta)}
Differentiate and apply boundary conditions.
\e{\fr{\p}{\p r}u(r,\theta)=\sum_{n=1}^\infty B_n 4n r^{4n-1} \sin(4n\theta)}
\e{\fr{\p}{\p r}u(1,\theta)=\sum_{n=1}^\infty B_n 4n \sin(4n\theta)=\sin(\theta)}
This is a Fourier sine series for $\sin(\theta)$. Expand the Fourier sine series, and solve for $B_n$.
\e{4nB_n &= \fr{2}{\theta_0} \int_0^{\theta_0} \sin(\theta) \sin(4n\theta)d\theta\\
&= \fr{8}{\pi}\int_0^{\fr{\pi}{4}} \fr{1}{2} (\cos((1-4n)\theta) - \cos((1+4n)\theta)) d\theta\\
&= \fr{4}{\pi} ((\fr{\sin((1-4n)\theta)}{1-4n})_0^{\fr{\pi}{4}}-(\fr{\sin((1+4n)\theta)}{1+4n})_0^{\fr{\pi}{4}})\\
&= \fr{4}{\pi} (\fr{(1+4n)\sin(\fr{\pi}{4}-n\pi)-(1-4n)\sin(\fr{\pi}{4}+n\pi)}{(1-4n)(1+4n)})\\
&= \fr{4}{\pi (1-16n^2)}((1+4n)\sin(\fr{\pi}{4})\cos(n\pi) - \cos(\fr{\pi}{4})\sin(n\pi))\\
&-(1-4n)(\sin(\fr{\pi}{4})\cos(n\pi) +\cos(\fr{\pi}{4})\sin(n\pi))\\
&= \fr{4}{\pi (1-16n^2)} ((1+4n)2^{-1/2} (-1)^n - (1-4n)2^{-1/2}(-1)^n)\\
&= \fr{4(-1)^n}{\pi(1-16n^2)2^{1/2}}(8n)\\
&= \fr{32(-1)^n n}{\pi(1-16n^2)2^{1/2}}}
Divide by $4n$ to get $B_n$.
\e{\boxed{B_n=\fr{8(-1)^n}{\pi(1-16n^2)2^{1/2}}}}
Apply $B_n$ to the solution for $u(r,\theta)$.
\e{\boxed{u(r,\theta)=\sum_{n=1}^\infty \fr{8(-1)^n}{\pi(1-16n^2)2^{1/2}} r^{4n} \sin(4n\theta)}}
\end{document}
